% ECONOMICS_PROBLEM: {"id": "H24-384", "title": "Wealth taxation", "jel_code": "H24", "source_id": "OP-0042"}
% !TeX program = xelatex
\documentclass[11pt,a4paper]{article}
\usepackage[margin=23mm]{geometry}
\usepackage{fontspec}
\setmainfont{Latin Modern Roman}
\usepackage{amsmath,amssymb,mathtools}
\usepackage{xcolor,enumitem,booktabs,longtable,array,xurl,hyperref}
\definecolor{muted}{HTML}{53636C}
\hypersetup{colorlinks=true,urlcolor=blue,linkcolor=blue}
\setlength{\parindent}{0pt}
\setlength{\parskip}{5pt}
\newcommand{\R}{\mathbb R}
\newcommand{\E}{\mathbb E}
\newcommand{\N}{\mathbb N}
\newcommand{\ind}{\mathbf 1}
\newcommand{\Var}{\operatorname{Var}}
\newcommand{\Cov}{\operatorname{Cov}}
\newcommand{\supp}{\operatorname{supp}}
\DeclareMathOperator*{\argmin}{arg\,min}
\DeclareMathOperator*{\argmax}{arg\,max}
\newcommand{\entrymeta}[3]{{\sffamily\small\color{muted}\textbf{\texttt{#1}}\quad|\quad #2\par #3\par}}
\newcommand{\Problemref}[1]{\texttt{#1}}
\newcommand{\JELref}[2]{\texttt{#2} (JEL \texttt{#1})}
\begin{document}
\section*{Wealth taxation}
\textbf{Problem ID:} \texttt{H24-384}\quad \textbf{JEL:} \texttt{H24}\par
% BEGIN ECONOMICS STATEMENT
\label{problem:OP-0042}
\label{atlas:prob:P2}
\noindent\textbf{Original atlas identifier:} P2; entry 42. \textbf{Field:} Public finance and taxation. \textbf{Original tractability label:} Hard.

\noindent\textbf{Classification:} Parameterized theoretical/empirical research agenda; not a certified unsolved theorem.

\subsection*{Definitions and assumptions}
Fix dates $t=0,\ldots,T$, discount factor $\beta\in(0,1]$, and an initial distribution $F$ over household wealth portfolios, liquidity, entrepreneurial ability, and residence. A recurrent-wealth policy $p$ specifies asset-specific tax bases, exemption thresholds, rates, valuation rules, reporting obligations, and enforcement. Economic wealth $A_{it}$, assessed wealth $\widehat A_{it}$, and reported wealth $\widetilde A_{it}$ are distinct variables for household $i$. Each model $m$ specifies investment, entrepreneurship, valuation error, avoidance, migration, and a selected general equilibrium. Let $EV_i(p;m)$ be discounted monetary equivalent variation, including private compliance and relocation costs. Let $R_t(p;m)$ be revenue across all affected taxes and $C_t(p;m)$ administrative resource costs. Fix nonnegative welfare weights $w_i$ and public-fund value $\lambda$; impose finite discounted moments.

\subsection*{Formal research question}
\emph{Editorial formulation: the mathematical target below is not a theorem quoted from the references.}

For a legally feasible class $\mathcal P$ satisfying enforcement-capacity limits and a stated present-value public budget, identify the welfare--revenue frontier
\[
 \mathcal F_m=\left\{\left(\sum_{t=0}^{T}\beta^t[R_t(p;m)-C_t(p;m)],\ \mathbb E_m[w_iEV_i(p;m)]\right):p\in\mathcal P\right\}.
\]
For each required net-revenue target $\bar R$, compare the supremum of weighted equivalent variation over policies meeting that target. Identify which asset-specific elasticities can be transported across valuation and reporting regimes, and give bounds when investment, reassessment, concealment, and mobility cannot be separately recovered. A policy dominates another only relative to the same welfare weights, horizon, equilibrium selection, and revenue requirement; a universal positive answer is not asserted.

\subsection*{Interpretation and scope}
“Severe” distortion must be replaced by an explicit welfare-loss or entrepreneurial-entry constraint chosen before estimation. Illiquid business wealth requires both valuation and financing rules; a low assessed value can reflect an exemption or measurement error rather than low economic wealth. Study migrants in origin and destination records to distinguish an actual international relocation from local tax-base reassignment. Specify whether entrepreneurs can use instalments, loss offsets, or deferrals, since those rules alter the borrowing constraint and investment response. Danish and Swiss institutional evidence can discipline a model, but extrapolation to a new tax base or national policy requires invariance assumptions. A useful empirical deliverable jointly measures net receipts, private burdens, asset prices, firm entry, and residence rather than treating reported-wealth growth as saving.

\subsection*{Source audit and status}
[S1] is an explicitly selected resolution of the ambiguous Saez--Zucman author-year citation. [S2] and [S3] provide existing country-specific evidence. None certifies that every design on this frontier is unstudied or that a globally optimal design is currently unknown.

\subsection*{References}
\begin{enumerate}[label={[S\arabic*]},leftmargin=*,itemsep=0.6em]
\item Emmanuel Saez and Gabriel Zucman (2019). Progressive Wealth Taxation. Brookings Papers on Economic Activity, Fall, 437--533.
\url{https://www.brookings.edu/articles/progressive-wealth-taxation/}
\newline\emph{Locator and support limits:} Paper and abstract: wealth-tax design and simulations, not direct identification of all long-run responses.
\item Katrine Jakobsen, Kristian Jakobsen, Henrik Kleven, and Gabriel Zucman (2020). Wealth Taxation and Wealth Accumulation: Theory and Evidence from Denmark. The Quarterly Journal of Economics 135(1), 329--388. DOI: 10.1093/qje/qjz032.
\url{https://academic.oup.com/qje/article-abstract/135/1/329/5584349}
\newline\emph{Locator and support limits:} Abstract and article: Danish wealth-tax reforms and accumulation; institutional transport requires additional assumptions.
\item Marius Brülhart, Jonathan Gruber, Matthias Krapf, and Kurt Schmidheiny (2022). Behavioral Responses to Wealth Taxes: Evidence from Switzerland. American Economic Journal: Economic Policy 14(4), 111--150. DOI: 10.1257/pol.20200258.
\url{https://www.aeaweb.org/articles?id=10.1257/pol.20200258}
\newline\emph{Locator and support limits:} Abstract: reported wealth, mobility, and housing-price responses in Swiss cantons; reported wealth is not real saving.
\end{enumerate}

\subsection*{Common conventions: Source mathematical conventions}

Unless an entry states otherwise, agent and firm sets are finite. State and action spaces are measurable, strategies and outcome maps are measurable, probability laws are specified, and expectations exist for the targets under discussion. Infinite-horizon discount factors lie in $(0,1)$ and discounted objectives are finite. Norms, tolerances and complexity input sizes must be specified for computational comparisons. Constants or primitives that appear locally are inputs, not empirical facts asserted by this catalogue.

\textbf{Identification.} A statistical model $m\in\mathcal M$ implies an observable law $P_m$ and target $\theta(m)$. At law $P$, its identified set is
\[
\Theta(P;\mathcal M)=\{\theta(m):m\in\mathcal M,\ P_m=P\}.
\]
Point identification holds when this set is a singleton, or equivalently when $P_{m_1}=P_{m_2}$ implies $\theta(m_1)=\theta(m_2)$ for all compatible models. Sharp bounds mean bounds attained, or approached when the boundary is not attained, by compatible models. Writing this set is a definition, not a solution: the substantive task is to establish informative restrictions and derive or estimate the set. An unrestricted-confounding model may give only trivial bounds.

\textbf{Inference.} Identification, estimability and valid uncertainty quantification are separate questions. Fix $\alpha\in(0,1)$. A confidence procedure $C_n$ over a declared law class $\mathcal P$ has asymptotic uniform coverage at level $1-\alpha$ when
\[
\liminf_{n\to\infty}\inf_{P\in\mathcal P}\Pr_P\{\theta(P)\in C_n\}\geq1-\alpha.
\]
For set-identified targets the coverage event must be defined separately, for example coverage of the full identified set. If an entry uses identification-indexed models instead of uniquely indexed laws, the same coverage convention is applied over those models. Pointwise limits do not establish uniform validity, and asymptotic validity does not guarantee finite-sample precision. Sample sizes, dependence structures and weak-identification sequences are part of each application.

\textbf{Counterfactuals.} $Y_i(a)$ is a potential outcome under a specified intervention, including its timing, implementation and versions. It is not $Y_i$ conditional on voluntarily choosing $a$. A full intervention vector permits interference; shorthand scalar treatment notation does not silently impose no interference. Assignment, selection, attrition, anticipation and measurement must be specified. A claim about an instrument's local effect is not automatically a population policy effect.

\textbf{Economic equilibrium.} A counterfactual model includes best responses, constraints, feasibility, market clearing, state transitions and expectations consistent with the stated information. When several equilibria exist, either a declared selector is an input or the target is set-valued. Comparisons across policy regimes must use the stated selection convention. Reduced-form relationships are not presumed invariant after an equilibrium-changing intervention.

\textbf{Optimization and welfare.} Optimization is over the stated feasible set. If attainment is not established, $\sup$ or $\inf$ and $\varepsilon$-optimal choices replace an asserted maximizer or minimizer. For example, with value $V=\sup_{a\in\mathcal A}W(a)<\infty$, an $\varepsilon$-optimal set is $\{a\in\mathcal A:W(a)\geq V-\varepsilon\}$ for $\varepsilon>0$. Benefits, costs, risk and health or environmental outcomes can be aggregated only using declared units and conversion weights. Interpersonal utility weights, the treatment of fiscal transfers and foreign profits, discounting, and the behavioral welfare criterion are explicit inputs. A comparative-static derivative additionally requires existence and appropriate differentiability of the selected counterfactual.

\textbf{Sources.} Local labels [S1], [S2], and so forth restart in every question. Locators identify the relevant abstract, model, section or result at the level actually checked; a bibliographic or abstract-level verification is not represented as inspection of an entire proof. Working papers are distinguished from later publications and changing versions. Source summaries are paraphrased. All URL checks and editorial formulations refer to the audit date stated above.
% END ECONOMICS STATEMENT
\end{document}
